\chapter{数列}\label{ch:b1:seq}

在高中分册中，数列已被运用，而极限概念
很大程度上靠信念接受。这里，理论重新建立在
$\R$ 的完备性（\cref{ch:b1:reals}）之上：每一条经典定理
--- 单调收敛、相邻数列、波尔查诺--魏尔斯特拉斯、
柯西准则 --- 都是这同一条公理的不同侧面。本章最后以
对由 $u_{n+1} =
f(u_n)$ 定义的数列的实际研究收尾。

\section{收敛}

\begin{definition}[数列的极限]\label{def:b1:seq:limit}
若实数列 $(u_n)$ 对 $\ell \in \R$ 满足
下式，则称其\emph{收敛}于该数\index{极限！数列的}：
\[
\forall \varepsilon > 0,\ \exists N \in \N,\ \forall n \geq N,
\qquad \abs{u_n - \ell} \leq \varepsilon .
\]
记作 $u_n \to \ell$ 或 $\lim u_n = \ell$。不收敛于任何实数的数列
称为\emph{发散}。\emph{趋于
$+\infty$}：$\forall M,\ \exists N,\ \forall n \geq N,\ u_n \geq M$
（$-\infty$ 的情形类似）。
\end{definition}

\begin{example}[完整写出一次的 $\varepsilon$--$N$ 证明]\label{ex:b1:seq:epsilonN}
断言：$u_n = \dfrac{n^2 + 1}{2n^2 - 3} \to \dfrac12$。先
把误差分离出来：
\[
\Bigl| u_n - \frac12 \Bigr|
= \Bigl| \frac{2(n^2 + 1) - (2n^2 - 3)}{2(2n^2 - 3)} \Bigr|
= \frac{5}{2\,\abs{2n^2 - 3}}
= \frac{5}{2\,(2n^2 - 3)} \quad (n \geq 2).
\]
然后用某个简单的式子从上方控制它：对 $n \geq 2$，$2n^2 - 3
\geq n^2$，故误差 $\leq \frac{5}{2n^2} \leq \frac 5{2n}$。
给定 $\varepsilon > 0$，由阿基米德性质可取 $N \geq
\max\bigl(2, \frac{5}{2\varepsilon}\bigr)$，则当 $n \geq N$ 时
误差 $\leq \varepsilon$。证毕。小结：一个
$\varepsilon$--$N$ 证明恰有三个步骤 --- 计算
误差，用一个递减的初等式给它定界，解出
阈值 --- 而在本章的定理（极限的运算法则、
夹逼准则）之后，几乎不会再写这样的证明：定理
已把三个步骤一次性打包好了。
\end{example}

\begin{example}[趋于无穷的发散，严格验证]\label{ex:b1:seq:toinfinity}
断言：$u_n = n^2 - 100n \to +\infty$。提出主导项：
$u_n = n^2\bigl(1 - \frac{100}{n}\bigr) \geq \frac{n^2}{2}$，当
$n \geq 200$ 时。给定 $M$，取 $N = \max\bigl(200,
\lceil\sqrt{2M}\rceil\bigr)$：当 $n \geq N$ 时，$u_n \geq
\frac{n^2}{2} \geq M$。本例体现两点做法：提出主导项
的做法把一场较量（$n^2$ 与 $-100n$）化为
单一尺度乘以趋于 $1$ 的因子；另一点是阈值
可以很大（$u_{100} = 0$，该数列甚至在
$n = 100$ 之前为负） --- 趋于 $+\infty$ 的发散是关于
尾部的断言，与任何有限次的反常行为无关。
\end{example}

\begin{proposition}[基本性质]\label{prop:b1:seq:first}
\begin{enumerate}
  \item 若极限存在，则它是唯一的。
  \item 收敛数列是有界的。
  \item 若 $u_n \to \ell$，则对有限个项的任意改动
        不改变数列的收敛性与极限。
\end{enumerate}
\end{proposition}

\begin{proof}
(1) 若 $u_n \to \ell$ 且 $u_n \to \ell'$，其中 $\ell \neq \ell'$，
取 $\varepsilon = \frac{\abs{\ell - \ell'}}{3}$，则超过两个
阈值时，$\abs{\ell - \ell'} \leq \abs{\ell - u_n} + \abs{u_n -
\ell'} \leq 2\varepsilon = \frac23 \abs{\ell - \ell'}$，矛盾。

(2) 取 $\varepsilon = 1$：超过 $N$ 后，$\abs{u_n} \leq \abs\ell +
1$；前面有限个项同样有界，因此 $\abs{u_n}
\leq \max(\abs{u_0}, \dots, \abs{u_{N-1}}, \abs\ell + 1)$。

(3) 详细地说：设 $v_n = u_n$ 对 $n \geq n_0$ 成立，且 $u_n \to
\ell$。给定 $\varepsilon > 0$，取阈值 $N$（对
$(u_n)$）：当 $n \geq \max(N, n_0)$ 时，$\abs{v_n - \ell} = \abs{u_n
- \ell} \leq \varepsilon$。于是 $v_n \to \ell$：定义
只对 $n \geq N$ 量化，任何有限的前段都可以
以放大阈值为代价被改写。（这就是为什么
“对所有充分大的 $n$”这类假设在
本章中处处都足够。）
\end{proof}

\begin{theorem}[极限的运算]\label{thm:b1:seq:operations}
若 $u_n \to \ell$ 且 $v_n \to m$，则
\[
u_n + v_n \to \ell + m, \qquad u_n v_n \to \ell m, \qquad
\frac{u_n}{v_n} \to \frac{\ell}{m} \ (\text{ 若 } m \neq 0),
\qquad \abs{u_n} \to \abs\ell .
\]
\end{theorem}

\begin{proof}
\emph{和：} $\abs{(u_n + v_n) - (\ell + m)} \leq \abs{u_n - \ell} +
\abs{v_n - m} \leq 2\varepsilon$，超过两个阈值中较大者后。
\emph{积：}写作
\[
u_n v_n - \ell m = (u_n - \ell)\,v_n + \ell\,(v_n - m);
\]
$(v_n)$ 有界，存在 $B$ 为其界（见\cref{prop:b1:seq:first}），故
右端 $\leq B\abs{u_n - \ell} + \abs{\ell}\,\abs{v_n - m}$，
可任意小。\emph{商：}只需处理 $\frac
1{v_n}$。取 $\varepsilon = \frac{\abs m}{2}$，则超过某个 $N_0$ 后，
$\abs{v_n} \geq \frac{\abs m}{2}$，于是
\[
\Bigl| \frac{1}{v_n} - \frac 1m \Bigr|
= \frac{\abs{m - v_n}}{\abs{v_n m}}
\leq \frac{2}{m^2}\,\abs{v_n - m} \longrightarrow 0 .
\]
\emph{绝对值：}$\bigl|\abs{u_n} - \abs\ell\bigr| \leq
\abs{u_n - \ell}$（反向三角不等式，
见\cref{prop:b1:complex:rules}）。
\end{proof}

\begin{example}[运算法则加一个代数技巧]\label{ex:b1:seq:conjugatetrick}
求 $\lim\,\bigl(\sqrt{n^2 + n} - n\bigr)$。两项
各自趋于 $+\infty$：运算定理对它们之差
不作任何断言（\emph{未定式}）。
乘以共轭式：
\[
\sqrt{n^2 + n} - n
= \frac{(n^2 + n) - n^2}{\sqrt{n^2 + n} + n}
= \frac{n}{\sqrt{n^2+n} + n}
= \frac{1}{\sqrt{1 + \frac1n} + 1} .
\]
现在各项都收敛：$\sqrt{1 + \frac1n} \to 1$，因为
$0 \leq \sqrt{1 + h} - 1 = \frac{h}{\sqrt{1+h} + 1} \leq h$
（再次共轭，再用 $h = \frac1n$ 夹逼）；随后
由运算定理得极限 $\frac{1}{1 + 1} =
\frac12$。要点是：运算定理并非
求所有极限的万能工具 --- 未定式（$\infty -
\infty$、$\frac00$、$0 \times \infty$、$1^\infty$）须先
经代数\emph{变形}（共轭、提取主导项），
直至各部分都收敛；对于顽固的情形，系统化的工具是
渐近展开，见
\cref{ch:b1:taylor}。
\end{example}

\begin{theorem}[极限与序关系]\label{thm:b1:seq:order}
\begin{enumerate}
  \item 若 $u_n \leq v_n$ 对所有充分大的 $n$ 成立，且二者都收敛，则
        $\lim u_n \leq \lim v_n$。（严格不等式\emph{不}
        能保持到极限：$\frac 1n > 0$ 而 $\lim = 0$。）
  \item （夹逼定理\index{夹逼定理}）若 $u_n \leq w_n
        \leq v_n$ 对所有充分大的 $n$ 成立，且 $u_n, v_n \to \ell$，则
        $w_n \to \ell$。
  \item 若 $u_n \to \ell > 0$，则 $u_n > \frac\ell2 > 0$ 对所有充分
        大的 $n$ 成立。
\end{enumerate}
\end{theorem}

\begin{proof}
（1）设 $\ell = \lim u_n > m = \lim v_n$；取 $\varepsilon =
\frac{\ell - m}{3}$，则充分大的项满足 $v_n \leq m + \varepsilon <
\ell - \varepsilon \leq u_n$，与 $u_n \leq v_n$ 矛盾。

（2）超过各自阈值后：$\ell - \varepsilon \leq u_n \leq w_n \leq
v_n \leq \ell + \varepsilon$。

（3）即\cref{def:b1:seq:limit}中取 $\varepsilon = \frac\ell2$ 的情形。
\end{proof}

\begin{example}[两次夹逼]\label{ex:b1:seq:squeezes}
（i）$\dfrac{\sin n}{n} \to 0$：由 $-\frac1n \leq \frac{\sin
n}{n} \leq \frac1n$，两侧界都夹向 $0$ --- 无需
去理解那不规则的分子。（ii）$(2^n +
3^n)^{1/n} \to 3$：先把根内部夹住，
\[
3^n \leq 2^n + 3^n \leq 2\cdot3^n
\quad\Longrightarrow\quad
3 \leq (2^n + 3^n)^{1/n} \leq 3\cdot 2^{1/n} ,
\]
而 $2^{1/n} = \eu^{\frac{\ln 2}{n}} \to 1$（与 $5^{1/n}$ 的情形
相同，见\cref{exo:b1:seq:2}），夹逼即得 $3$。要点在于
竞争指数之和的行为类似于其
\emph{最大}项 --- 较小的各项被一个无害的常数因子吸收，
这一常数因子随后被 $n$ 次方根消去。
\end{example}

\section{单调数列}

\begin{theorem}[单调极限定理]\label{thm:b1:seq:monotone}
有上界的递增数列收敛，极限为 $\sup\{u_n : n \in
\N\}$；无上界的递增数列趋于
$+\infty$。（对递减数列
有对称的结论。）\index{单调极限定理}
\end{theorem}

\begin{proof}
设 $s = \sup\{u_n\}$（见\cref{thm:b1:reals:sup}）。给定 $\varepsilon >
0$，由上确界的 $\varepsilon$ 刻画
（见\cref{prop:b1:reals:epsilon}）可得 $N$，使 $u_N > s -
\varepsilon$；由单调性，$s - \varepsilon < u_N \leq u_n \leq
s$（对所有 $n \geq N$），即收敛于 $s$。若无界：对每个
$M$，存在 $u_N > M$，单调性保证此后各项都大于 $M$。
\end{proof}

\begin{example}[作为存在性工具的单调定理]\label{ex:b1:seq:existencemachine}
设 $u_n = \prod_{k=1}^{n} \bigl(1 + \frac{1}{2^k}\bigr)$。每个
因子都大于 $1$，故 $(u_n)$ 递增。是否有上界？取
对数，并利用 $\ln(1 + x) \leq x$
（\cref{ex:b1:derivative:convexineq} 已预先给出；或用
高中分册中的粗略不等式 $1 + x \leq \eu^x$）：
\[
\ln u_n = \sum_{k=1}^{n} \ln\Bigl(1 + \frac{1}{2^k}\Bigr)
\leq \sum_{k=1}^{n} \frac{1}{2^k} < 1 ,
\]
故 $u_n < \eu$。递增且有界，$(u_n)$ 收敛于某个
$\ell \in \intoc{u_1}{\eu}$ --- 一个完全良定义的实数，
看不出闭式（$\ell = 2.384\dots$）。要点在于
单调极限定理是分析中最廉价的存在性工具；
它给出了 $\eu$ 本身
（见下文\cref{ex:b1:seq:e}），而在\cref{ch:b1:series}中将
仅凭
有界性判定每个正项级数的收敛性。
\end{example}

\begin{theorem}[相邻数列]\label{thm:b1:seq:adjacent}
设 $(a_n)$ 递增，$(b_n)$ 递减，且 $b_n - a_n \to
0$。则二者都收敛于同一个\emph{公共}极限 $\ell$，且 $a_n
\leq \ell \leq b_n$ 对所有 $n$ 成立。\index{相邻数列}
\end{theorem}

\begin{proof}
首先，$a_n \leq b_n$ 对所有 $n$ 成立：数列 $(b_n - a_n)$ 递减
且趋于 $0$，故其各项 $\geq 0$（若某项为负，则
将始终小于 $0$）。随后 $(a_n)$ 递增且以
$b_0$ 为上界，故收敛于某个 $\ell$（见\cref{thm:b1:seq:monotone}）；
同理 $(b_n) \to \ell'$；且 $\ell' - \ell = \lim (b_n - a_n) =
0$。不等式 $a_n \leq \ell \leq b_n$ 由
单调性推出（$\ell = \sup a_k \geq a_n$ 等）。
\end{proof}

\begin{omfigure}
\begin{tikzpicture}[scale=10.5]
  \draw[-Stealth] (-0.03,0) -- (1.08,0);
  % a_n increasing (omDef), b_n decreasing (omProp)
  \fill[omDef] (0,0) circle (0.3pt) node[below=2pt, font=\small]
    {$a_0$};
  \fill[omDef] (0.18,0) circle (0.3pt) node[below=2pt,
    font=\small] {$a_1$};
  \fill[omDef] (0.30,0) circle (0.3pt) node[below=2pt,
    font=\small] {$a_2$};
  \fill[omDef] (0.37,0) circle (0.3pt);
  \fill[omDef] (0.41,0) circle (0.3pt);
  \fill[omProp] (1.0,0) circle (0.3pt) node[below=2pt,
    font=\small] {$b_0$};
  \fill[omProp] (0.78,0) circle (0.3pt) node[below=2pt,
    font=\small] {$b_1$};
  \fill[omProp] (0.62,0) circle (0.3pt) node[below=2pt,
    font=\small] {$b_2$};
  \fill[omProp] (0.54,0) circle (0.3pt);
  \fill[omProp] (0.50,0) circle (0.3pt);
  \draw[omThm, thick] (0.455,-0.014) -- (0.455,0.014);
  \node[above, omThm, font=\small] at (0.455,0.02) {$\ell$};
  \draw[Stealth-Stealth, gray] (0.30,0.05) -- (0.62,0.05);
  \node[above, font=\small, text=gray] at (0.46,0.055)
    {$b_n - a_n \to 0$};
\end{tikzpicture}

{\small 相邻数列：$(a_n)$ 递增，$(b_n)$ 递减，
两者之间的间距趋于 $0$。每个区间
$\intcc{a_n}{b_n}$ 包含之后的所有区间，而公共极限
$\ell$ 是唯一留在每个区间中的点 --- 这幅图背后是
下文波尔查诺--魏尔斯特拉斯定理的二分法证明，
以及在
\cref{ch:b1:continuity}中的介值定理的直观背景。}
\end{omfigure}

\begin{example}[自然常数 $\eu$]\label{ex:b1:seq:e}
设 $a_n = \sum_{k=0}^{n} \frac{1}{k!}$，$b_n = a_n + \frac{1}{n
\cdot n!}$（$n \geq 1$）。则 $(a_n)$ 递增；且
\[
b_{n+1} - b_n = \frac{1}{(n+1)!} + \frac{1}{(n+1)(n+1)!} -
\frac{1}{n\,n!}
= \frac{n(n+1) + n - (n+1)^2}{n(n+1)(n+1)!}
= \frac{-1}{n(n+1)(n+1)!} < 0 ,
\]
于是 $(b_n)$ 递减，且 $b_n - a_n \to 0$，故二者相邻。它们的公共
极限（此处按定义）即为自然常数 $\eu \approx 2.71828$；
不等式 $a_n < \eu < b_n$ 足以证明 $\eu \notin
\Q$（见\cref{exo:b1:seq:9}）。
\end{example}

\section{子列与波尔查诺--魏尔斯特拉斯定理}

\begin{definition}[子列]\label{def:b1:seq:subsequence}
称 $(u_n)$ 的\emph{子列}\index{子列}为
数列 $(u_{\varphi(n)})$，其中 $\varphi \colon \N \to \N$ 严格递增
（由归纳法可得 $\varphi(n) \geq n$）。
\end{definition}

\begin{proposition}\label{prop:b1:seq:subsequences}
若 $u_n \to \ell$（$\ell \in \R$ 或 $\pm\infty$），则其每个子列
都趋于 $\ell$。由此，若数列有两个
极限不同的子列，则该数列发散。反之，若 $(u_{2n})$ 与
$(u_{2n+1})$ 都收敛于\emph{同一个} $\ell$，则 $u_n \to
\ell$。
\end{proposition}

\begin{proof}
超过阈值 $N$（对 $(u_n)$ 而言）之后，所有指标 $\varphi(n) \geq n
\geq N$ 均满足要求（不等式 $\varphi(n) \geq n$ 正是
\cref{def:b1:seq:subsequence}中所述的归纳：$\varphi(0)
\geq 0$，且 $\varphi(n+1) > \varphi(n) \geq n$ 推出
$\varphi(n+1) \geq n + 1$）。对于逆命题，给定
$\varepsilon$，取两个阈值 $N_0$（偶数项）和 $N_1$
（奇数项）；任取指标 $n \geq \max(2N_0, 2N_1 + 1)$，则
其或为偶数，$n = 2k$ 且 $k \geq N_0$；或为奇数，$n = 2k+1$ 且
$k \geq N_1$ --- 两种情形下均有 $\abs{u_n - \ell} \leq
\varepsilon$：每个指标都被两个子列之一覆盖，
这正是证明的要点。
\end{proof}

\begin{example}[子列极限]\label{ex:b1:seq:sublimits}
对于 $u_n = (-1)^n \frac{n}{n+1}$：偶数项构成的子列趋于
$1$，奇数项子列趋于 $-1$，故该数列发散 --- 但其
发散方式是有序的，聚集在两个值 $\pm 1$ 附近。
对于 $u_n = \cos\frac{2\pi n}{3}$：取下标分别为
$3k$、$3k + 1$、$3k + 2$ 的三个子列均为常数，分别等于 $1$、
$-\frac12$、$-\frac12$；子列极限的集合为 $\{1,
-\frac12\}$。总结要点是：\emph{有界}数列
当且仅当它恰有一个子列极限时收敛
（见\cref{exo:b1:seq:8}）；有界数列的发散总是
意味着至少有两个聚点，而下文的波尔查诺--魏尔斯特拉斯定理
保证至少存在一个子列极限。
\end{example}

\begin{theorem}[波尔查诺--魏尔斯特拉斯定理]\label{thm:b1:seq:bw}
每个有界实数列都有一个
收敛子列。\index{波尔查诺--魏尔斯特拉斯定理}
\end{theorem}

\begin{proof}
设 $u_n \in \intcc{a}{b}$（对所有 $n$）。用二分法构造嵌套
区间：取 $\intcc{a_0}{b_0} = \intcc{a}{b}$；给定
$\intcc{a_k}{b_k}$ 包含 $u_n$（对无穷多个 $n$），则其两半之一
仍包含 $u_n$（对无穷多个 $n$） --- 记为
$\intcc{a_{k+1}}{b_{k+1}}$。数列 $(a_k)$、$(b_k)$ 相邻
（$b_k - a_k = \frac{b-a}{2^k} \to 0$），其公共极限为
$\ell$（见\cref{thm:b1:seq:adjacent}）。

取子列：选取 $\varphi(0)$，使 $u_{\varphi(0)} \in
\intcc{a_0}{b_0}$；再归纳地，取 $\varphi(k+1) > \varphi(k)$
使 $u_{\varphi(k+1)} \in \intcc{a_{k+1}}{b_{k+1}}$ --- 这是可能的
因为该区间含有无穷多项。于是 $a_k \leq
u_{\varphi(k)} \leq b_k$，由夹逼定理得
$u_{\varphi(k)} \to \ell$。
\end{proof}

\begin{remark}[波尔查诺--魏尔斯特拉斯定理说了什么，没说什么]
它\emph{确实}断言：仅由有界性，便有某个子列收敛 ---
存在性而不给出公式，正如二分法证明所清楚表明的
（没有任何东西告诉我们\emph{哪些}指标被保留）。它
\emph{并不}断言极限唯一：数列 $((-1)^n)$ 的子列分别收敛于 $1$ 与 $-1$，
而子列极限构成的集合
甚至可以是无穷集
（\cref{ex:b1:seq:sublimits}，以及
\cref{pb:b1:topology:1}中的整个康托集）。它不能推广到无界情形：
数列 $(n)$ 完全没有收敛子列 --- 不过，
对任何无界数列，总能抽取一个趋于 $+\infty$ 或 $-\infty$ 的子列
（例如取 $\varphi(k)$，使
$u_{\varphi(k)} \geq k$）。用得好时，该定理是一台
\emph{存在性发生器}：它出现在下文柯西准则的关键处，
也出现在海涅定理与极值定理中 ---
总是为了得到一个显式构造所不能给出的点，
这正是它的价值所在。
\end{remark}

\section{柯西序列与完备性}

\begin{definition}\label{def:b1:seq:cauchy}
称数列 $(u_n)$ 为\emph{柯西序列}\index{柯西序列}
当其各项\emph{彼此}任意接近时：
\[
\forall \varepsilon > 0,\ \exists N,\ \forall p, q \geq N,
\qquad \abs{u_p - u_q} \leq \varepsilon .
\]
\end{definition}

\begin{example}[用手算验证柯西性质]\label{ex:b1:seq:cauchyhand}
设 $u_n = \sum_{k=0}^{n} \frac{\cos k}{2^k}$------既不单调，
也没有可猜的极限。对 $p > q$：
\[
\abs{u_p - u_q}
= \Bigl| \sum_{k=q+1}^{p} \frac{\cos k}{2^k} \Bigr|
\leq \sum_{k=q+1}^{p} \frac{1}{2^k}
< \frac{1}{2^{q}} ,
\]
由三角不等式、$\abs{\cos k} \leq 1$ 及有限的
等比和。任给 $\varepsilon > 0$，取 $N$ 使 $2^{-N}
\leq \varepsilon$：则 $N$ 之后的所有间距都 $\leq \varepsilon$，
于是数列为柯西序列，从而收敛 --- 极限
无法写成闭形式，这正是关键所在。本例的要点是
以几何级数控制增量，这是证明柯西性质的标准做法，
在\cref{ch:b1:series}中，这一论证将被归结为
“绝对收敛蕴含收敛
”。
\end{example}

\begin{theorem}[$\R$ 的完备性]\label{thm:b1:seq:complete}
实数列收敛，当且仅当它是
柯西序列。\index{$\R$ 的完备性}
\end{theorem}

\begin{proof}
（$\Rightarrow$）若 $u_n \to \ell$，则当
$\frac\varepsilon2$ 对应的阈值之后，$\abs{u_p - u_q} \leq \abs{u_p - \ell} +
\abs{\ell - u_q} \leq \varepsilon$.

（$\Leftarrow$）设 $(u_n)$ 为柯西序列。\emph{它有界}：取
$\varepsilon = 1$，则 $N$ 之后各项都落在半径为 $1$、以 $u_N$ 为中心的邻域内，
而前面的有限个项有界，从而数列有界。\emph{抽取子列}：由
\cref{thm:b1:seq:bw}，存在子列 $u_{\varphi(n)} \to \ell$。
\emph{结论}：任给 $\varepsilon > 0$，取 $N$（由柯西性质，对应于
$\frac\varepsilon2$），并取 $n \geq N$，使 $\abs{u_{\varphi(n)} - \ell}
\leq \frac\varepsilon2$ 且 $\varphi(n) \geq N$；于是对每个 $p
\geq N$，有
\[
\abs{u_p - \ell} \leq \abs{u_p - u_{\varphi(n)}} +
\abs{u_{\varphi(n)} - \ell} \leq \varepsilon . \qedhere
\]
\end{proof}

\begin{remark}
该准则的价值在于：它\emph{无需指出极限
}便可证明收敛。它在 $\Q$ 上不成立（十进小数截断
$\sqrt 2$ 的截断构成有理数柯西序列，却无有理极限）：
完备性是 $\R$ 的性质，等价于上
确界公理。它也是级数收敛背后的主要工具
（见\cref{ch:b1:series}）。
\end{remark}

\begin{example}[极限不可见的柯西序列]\label{ex:b1:seq:basel}
设 $S_n = \sum_{k=1}^{n} \frac{1}{k^2}$。对 $p > q \geq 1$：
\[
S_p - S_q = \sum_{k=q+1}^{p} \frac{1}{k^2}
\leq \sum_{k=q+1}^{p} \frac{1}{k(k-1)}
= \sum_{k=q+1}^{p} \Bigl(\frac{1}{k-1} - \frac 1k\Bigr)
= \frac 1q - \frac 1p < \frac 1q ,
\]
于是取 $N > \frac1\varepsilon$，则此后所有间距都 $\leq \varepsilon$：
$(S_n)$ 为柯西序列，从而收敛。注意刚才的论证：
证明了某个具体实数存在，却未给它任何名称，
（它就是 $\frac{\pi^2}{6}$ --- 一个著名的
欧拉恒等式，已在本科二年级分册中证明；本章中没有任何内容能
说明这一点）。这种分工 --- 先证存在，
之后（若有必要）再确定其值 --- 正是柯西准则的全部要旨，
也是级数理论的动力，
在\cref{ch:b1:series}中。
\end{example}

\section{递推数列}

\begin{method}[研究 $u_{n+1} = f(u_n)$]\label{met:b1:seq:recurrent}
给定 $f$ 及初值 $u_0$：\index{递推数列}
\begin{enumerate}
  \item \emph{稳定区间}：求区间 $I$，使 $f(I)
        \subseteq I$ 且包含 $u_0$，则所有 $u_n \in I$（由
        归纳法）。
  \item \emph{候选极限}：若 $u_n \to \ell \in I$，且 $f$
        在 $\ell$ 处连续（见\cref{ch:b1:continuity}），则 $\ell$
        为\emph{不动点}：$f(\ell) = \ell$。求解 $f(x) = x$。
  \item \emph{单调性}：若 $f$ 在 $I$ 上递增，则
        $(u_n)$ 单调（若 $u_1 \geq u_0$ 则递增，
        否则递减）；结合有界性，
        由\cref{thm:b1:seq:monotone}即得结论。若 $f$ 递减，
        则考察两个子列 $(u_{2n})$ 与 $(u_{2n+1})$，
        它们对 $f \circ f$ 而言是单调的。
  \item \emph{误差控制}：若不等式 $\abs{f(x) - \ell} \leq
        k\abs{x - \ell}$（其中 $k < 1$），则 $\abs{u_n - \ell} \leq
        k^n \abs{u_0 - \ell} \to 0$，直接可得。
\end{enumerate}
\end{method}

\begin{example}[海伦方法]\label{ex:b1:seq:heron}
设 $u_0 = 2$，$u_{n+1} = \dfrac12\Bigl(u_n +
\dfrac{2}{u_n}\Bigr)$：求
$\sqrt 2$ 的古老算法。\index{海伦方法}
\begin{itemize}
  \item \emph{稳定性：}当 $x > 0$ 时，由算术--几何平均
        不等式给出 $\frac12(x + \frac2x) \geq \sqrt{x \cdot
        \frac 2x} = \sqrt 2$，故 $I = \intco{\sqrt 2}{+\infty}$ 是
        稳定的，且包含 $u_1$（事实上 $u_1 = \frac32 \geq
        \sqrt 2$）。
  \item \emph{单调性：}当 $x \geq \sqrt 2$ 时，$\;x - f(x) =
        \frac{x^2 - 2}{2x} \geq 0$：数列自
        $u_1$ 起递减，且以 $\sqrt 2$ 为下界，因而收敛。
  \item \emph{极限：}不动点满足 $x = \frac12(x +
        \frac2x)$，即 $x^2 = 2$；在 $I$ 上，$\ell = \sqrt 2$。
  \item \emph{收敛速度：}$u_{n+1} - \sqrt 2 = \frac{(u_n -
        \sqrt2)^2}{2u_n} \leq \frac{(u_n - \sqrt 2)^2}{2\sqrt 2}$：
        正确数字的位数在每一步大致 \emph{翻一番}
        （二次收敛）。
\end{itemize}
\end{example}

\begin{omfigure}
\begin{tikzpicture}
\begin{axis}[omaxis, width=8.6cm, height=7cm,
    xmin=0.9, xmax=2.35, ymin=0.9, ymax=2.35,
    xtick={1,1.5,2}, ytick={1,1.5,2},
    xlabel={$x$}, ylabel={$y$}]
  \addplot[gray, dashed, domain=0.9:2.3] {x}
    node[pos=0.72, above left, font=\small] {$y = x$};
  \addplot[omDef, very thick, domain=0.95:2.3, samples=100]
    {0.5*(x + 2/x)} node[pos=0.88, above left, font=\small]
    {$y = f(x)$};
  % cobweb: u0=2, u1=1.5, u2=1.4167
  \draw[gray, dashed, thin] (axis cs:1.5,0.9) -- (axis cs:1.5,1.4167);
  \draw[omThm, thick, -Stealth] (axis cs:2,0.9) -- (axis cs:2,1.5);
  \draw[omThm, thick, -Stealth] (axis cs:2,1.5) -- (axis cs:1.5,1.5);
  \draw[omThm, thick, -Stealth] (axis cs:1.5,1.5) --
    (axis cs:1.5,1.4167);
  \draw[omThm, thick, -Stealth] (axis cs:1.5,1.4167) --
    (axis cs:1.4167,1.4167);
  \fill (axis cs:1.41421,1.41421) circle (1.6pt);
  \node[anchor=north, font=\small] at (axis cs:1.4142,1.385)
    {$\sqrt 2$};
  \node[anchor=west, font=\small, omThm] at (axis cs:2.02,1.05)
    {$u_0$};
  \node[anchor=east, font=\small, omThm] at (axis cs:1.48,1.05)
    {$u_1$};
\end{axis}
\end{tikzpicture}

{\small 海伦迭代 $u_{n+1} = \frac12\bigl(u_n +
\frac{2}{u_n}\bigr)$，图中阶梯折线位于 $f$ 的图像
与对角线 $y = x$ 之间：由 $u_0 = 2$ 出发，迭代值沿阶梯
滑向不动点 $\sqrt 2$。}
\end{omfigure}

\begin{remark}[关于极限的常见误区]
四个经典例子。(i) \emph{步长趋于零不蕴含收敛}：
$u_{n+1} - u_n \to 0$ 远弱于柯西条件 ---
调和级数的部分和 $H_n$ 的步长 $\frac{1}{n+1} \to 0$，却
发散到 $+\infty$（见\cref{exo:b1:seq:5}）；柯西条件
控制充分大的\emph{所有}数对的 $\abs{u_p - u_q}$，而非
相邻的那些。(ii) \emph{严格不等式在
极限中失效}：由 $u_n < v_n$（对所有 $n$），只能得到 $\lim u_n
\leq \lim v_n$（见\cref{thm:b1:seq:order}）；$\frac1n > 0$，但
$\lim = 0$。(iii) \emph{有界不蕴含收敛}：$((-1)^n)$
是有界的却发散；有界加上 \emph{单调性}
则收敛；仅有界只能保证一个收敛的
子列（见\cref{thm:b1:seq:bw}）。(iv) \emph{不动点
排在第二位，而非第一位}：对 $u_{n+1} = f(u_n)$，
求解 $f(\ell) = \ell$ 用来确定极限，\emph{仅在}
收敛性证明之后。递推式 $u_{n+1} = 2u_n$ 的
唯一不动点为 $\ell = 0$，但由 $u_0 = 1$ 出发，数列
趋于 $+\infty$：方程 $\ell = 2\ell$ 从来
没有资格给出极限。运算顺序始终是：先证存在性
（见\cref{met:b1:seq:recurrent}，第 1--3 步），再确定
极限。
\end{remark}

\begin{example}[递减函数 $f$：黄金比递推式]\label{ex:b1:seq:golden}
设 $u_0 = 1$，$u_{n+1} = \dfrac{1}{1 + u_n}$。此处 $f(x) =
\frac{1}{1+x}$ 是 \emph{递减}的，故数列不单调
（它绕极限交替摆动），此时应采用
\cref{met:b1:seq:recurrent} 中的压缩步骤。稳定性：若
$x \in \intcc{\frac12}{1}$，则 $1 + x \in
\intcc{\frac32}{2}$，故 $f(x) \in \intcc{\frac12}{\frac23}
\subseteq \intcc{\frac12}{1}$，且 $u_1 = \frac12$，于是整个数列
都在其中。不动点：$\ell = \frac{1}{1+\ell}$
由 $\ell > 0$ 得 $\ell^2 + \ell - 1 = 0$，即
\[
\ell = \frac{\sqrt5 - 1}{2} = 0.6180\dots
\]
（黄金比的倒数）。压缩性：对 $x, y \in
\intcc{\frac12}{1}$，
\[
\abs{f(x) - f(y)} = \frac{\abs{x - y}}{(1+x)(1+y)}
\leq \frac{\abs{x-y}}{(3/2)^2} = \frac49\,\abs{x - y} ,
\]
故 $\abs{u_n - \ell} \leq \bigl(\frac49\bigr)^{n-1}\abs{u_1 -
\ell} \to 0$：由此收敛，速度呈几何级数型，
无需单调性。结语：单调方法与
压缩方法共同划分了递推数列的世界 ---
递增的 $f$ 给出单调轨道，递减的 $f$ 给出
被利普希茨常数 $< 1$ 驯服的交替轨道（
系统理论见\cref{exo:b1:derivative:11}）。
\end{example}

\begin{remark}[本分册中的展望]
数列是本分册其余部分的度量工具，可以衡量每一个对象。在
\cref{ch:b1:topology}中，数列
\emph{刻画}闭性与紧性；在
\cref{ch:b1:continuity}中，数列传递函数的极限；在
\cref{ch:b1:integration}中，黎曼和是收敛于积分的数列；
\cref{ch:b1:series}\emph{正是}一类特殊数列的理论，
即部分和。代数各章同样用到数列：
矩阵的迭代在
\cref{ch:b1:matrices}中构成数列，其性态
（$A^n$ 的收敛性）是一个线性代数问题，须用
本章的词汇来表述。需要处处牢记的两个定理：
单调极限定理（由序关系得存在性）与
波尔查诺--魏尔斯特拉斯定理（由有界得存在性） --- 二者合在一起，本书中几乎每一个
极限都由此诞生。
\end{remark}

\begin{remark}[复数列]
复数列 $(z_n)$ 收敛于 $\ell$，当
$\abs{z_n - \ell} \to 0$；等价地，当 $\Re(z_n) \to \Re(\ell)$
且 $\Im(z_n) \to \Im(\ell)$（比较 $\abs{z}$ 与 $\abs{\Re z} +
\abs{\Im z}$）。不涉及序的定理 --- 运算法则、
波尔查诺--魏尔斯特拉斯定理（两次取子列）、柯西收敛准则 --- 均可
原样移植。
\end{remark}

\section{练习}

\begin{exercise}[$\star$]\label{exo:b1:seq:1}
直接依据\cref{def:b1:seq:limit}，证明
$\dfrac{2n+1}{n+3} \to 2$，并证明 $(u_n) = ((-1)^n)$ 发散。
\end{exercise}

\begin{exercise}[$\star$]\label{exo:b1:seq:2}
求下列极限：
\[
\frac{n^2 - 3n + 1}{2n^2 + 5},
\qquad
\sqrt{n+1} - \sqrt n,
\qquad
\frac{2^n + n^3}{3^n - n^2},
\qquad
\sqrt[n]{5}\ \Bigl(= 5^{1/n}\Bigr).
\]
\end{exercise}

\begin{exercise}[$\star$]\label{exo:b1:seq:3}
证明标准的比较结论：若 $\abs{q} < 1$，则 $q^n \to 0$
\emph{（记 $\frac{1}{\abs q} = 1 + h$，$h > 0$，并利用
伯努利不等式 $(1+h)^n \geq 1 + nh$，其证明可用
归纳法）}。当 $q = 1$、$q = -1$、$\abs q >
1$ 时，数列的性态如何？
\end{exercise}

\begin{exercise}[$\star$]\label{exo:b1:seq:4}
设 $u_{n+1} = \frac{u_n + 3}{2}$，$u_0 = 0$。求不动点
$\ell$，证明 $v_n = u_n - \ell$ 为等比数列，并
求其显式通项公式及 $(u_n)$ 的极限。
\end{exercise}

\begin{exercise}[$\star\star$]\label{exo:b1:seq:5}
（调和级数）记 $H_n = \sum_{k=1}^{n} \frac 1k$，证明
$H_{2n} - H_n \geq \frac12$ 对所有 $n \geq 1$ 成立，并由此推出
$(H_n)$ \emph{不}是柯西序列，因而发散（趋于
$+\infty$，因其递增）。
\end{exercise}

\begin{exercise}[$\star\star$]\label{exo:b1:seq:6}
设 $(u_{2n})$、$(u_{2n+1})$ 与 $(u_{3n})$ 均收敛，证明
$(u_n)$ 收敛。\emph{（取公共子列，
以证明两者极限相等。）}
\end{exercise}

\begin{exercise}[$\star\star$]\label{exo:b1:seq:7}
研究数列 $u_0 = 0$，$u_{n+1} = \sqrt{2 + u_n}$：
其稳定性、单调性与极限。再证明误差估计
$\abs{u_n - 2} \leq \dfrac{2}{3^{\,n}}$ \emph{（证明 $2 - u_{n+1} =
\dfrac{2 - u_n}{2 + \sqrt{2 + u_n}}$，并将分母的下界取为
$3$）}。
\end{exercise}

\begin{exercise}[$\star\star$]\label{exo:b1:seq:8}
设 $(u_n)$ 有界，且 $(u_n)$ 的每个收敛子列
的极限都是\emph{同一个} $\ell$。证明 $u_n \to \ell$。
\emph{（用反证法及波尔查诺--魏尔斯特拉斯定理。）}
\end{exercise}

\begin{exercise}[$\star\star\star$]\label{exo:b1:seq:9}
沿用\cref{ex:b1:seq:e}中的记号，设 $\eu = \frac pq$
其中 $p, q \in \N^*$。利用 $a_q < \eu < b_q = a_q + \frac{1}{q\,
q!}$，两边乘以 $q!$，推出两整数之间的
矛盾。由此得 $\eu$ 为无理数。
\end{exercise}

\begin{exercise}[$\star\star\star$]\label{exo:b1:seq:10}
（切萨罗平均）对数列 $(u_n)_{n \geq 1}$，记 $c_n =
\frac{u_1 + \dots + u_n}{n}$。
\begin{enumerate}
  \item 证明：若 $u_n \to \ell$，则 $c_n \to \ell$\emph{（在阈值
        $N$ 处截断和式；前段以固定量
        除以 $n$ 估计，后段以 $\varepsilon$ 估计）}。
  \item 举例说明逆命题不成立。
  \item 推出：若 $u_{n+1} - u_n \to \ell$，则 $\frac{u_n}{n}
        \to \ell$。
\end{enumerate}
\end{exercise}

\begin{exercise}[$\star\star\star$]\label{exo:b1:seq:11}
设 $(u_n)$ 满足 $0 \leq u_{m+n} \leq u_m + u_n$，对所有 $m, n$
（次可加性）。证明 $\bigl(\frac{u_n}{n}\bigr)$ 收敛于
$\inf_{n \geq 1} \frac{u_n}{n}$。\emph{（固定 $m$，写 $n = qm
+ r$，估计 $\frac{u_n}{n}$ 时利用 $u_n \leq q\,u_m + u_r$。）}
\end{exercise}

\begin{exercise}[$\star\star\star$]\label{exo:b1:seq:12}
利用 $\Z + 2\pi\Z$ 的稠密性（它是 $(\R, +)$ 的子群）
（见\cref{exo:b1:reals:9}），证明数列 $(\sin n)_{n \in
\N}$ 在 $\intcc{-1}{1}$ 中稠密 --- 特别地，它发散。
\end{exercise}

\section{问题：切萨罗、施托尔茨与正弦的缓慢下降}

\begin{problem}[{周末问题 --- 切萨罗--施托尔茨定理，
以及 $u_n \sim \sqrt{3/n}$ 的渐近性，对于 $u_{n+1} = \sin
u_n$}]
\label{pb:b1:seq:1}
切萨罗--施托尔茨定理是离散形式的洛必达法则：欲求
商 $a_n/b_n$ 的极限，只需求
\emph{增量}之商 $(a_{n+1} -
a_n)/(b_{n+1} - b_n)$ 的极限。本题证明该定理，
借此求若干经典极限，再将其用于一个著名的例子：
数列 $u_{n+1} = \sin u_n$，它缓慢地趋于 $0$，
速度可精确计算，其渐近形式为 $u_n \sim \sqrt{3/n}$。\index{切萨罗--施托尔茨
定理} 高中分册中的两个事实在此承认，
并将在本分册后文中再予严格证明。其一是
切线不等式：
\[
\tag{G1} \eu^{u} \geq 1 + u \quad (u \in \R),
\]
在\cref{ch:b1:derivative}中借凸性重新证明，另一个是正弦的
夹逼估计：
\[
\tag{G2} x - \frac{x^3}{6} \;\leq\; \sin x \;\leq\; x -
\frac{x^3}{6} + \frac{x^5}{120} \quad (0 \leq x \leq 1), \qquad
\abs{\sin x} \leq \abs{x} \quad (x \in \R),
\]
在\cref{ch:b1:taylor}中借泰勒公式重新证明。

\medskip
\noindent\textbf{第一部分 --- 没有闭式的求和。}
\begin{enumerate}
  \item 利用 $1 + 2 + \dots + n = \frac{n(n+1)}{2}$ 与 $1^2 +
        \dots + n^2 = \frac{n(n+1)(2n+1)}{6}$，计算
        $\lim \frac{1 + 2 + \dots + n}{n^2}$ 与 $\lim \frac{1^2 +
        \dots + n^2}{n^3}$。
  \item 设 $T_n = \sum_{k=1}^n \sqrt k$，对此不存在闭式。
        证明下面的夹逼不等式：
        \[
        \frac{1}{2\sqrt 2}\,n^{3/2} \;\leq\; T_n \;\leq\;
        n^{3/2}
        \]
        \emph{（求下界时只保留满足 $k > \frac n2$ 的
        项）}。于是 $T_n$ 的阶为 $n^{3/2}$，
        但常数是多少？这一问题留到第 8 题再讨论。
  \item （裂项相消引理）设 $(b_n)$ 严格递增，
        且假设对所有 $k \geq N$，
        \[
        m \;\leq\; \frac{a_{k+1} - a_k}{b_{k+1} - b_k} \;\leq\;
        M .
        \]
        证明 $m \leq \dfrac{a_n - a_N}{b_n - b_N} \leq M$
        对每个 $n > N$ 成立。
\end{enumerate}

\medskip
\noindent\textbf{第二部分 --- 切萨罗--施托尔茨定理。}设
        $(b_n)$ 严格递增，且 $b_n \to +\infty$，
        并假设 $\dfrac{a_{n+1} - a_n}{b_{n+1} - b_n} \to \ell \in \R$。
\begin{enumerate}[resume]
  \item 任取 $\varepsilon > 0$，证明存在 $N$，使得
        $\ell - \varepsilon \leq \dfrac{a_n - a_N}{b_n - b_N}
        \leq \ell + \varepsilon$ 对所有 $n > N$ 成立。
  \item 对 $n > N$，证明恒等式
        \[
        \frac{a_n}{b_n} - \ell
        = \frac{a_N - \ell\,b_N}{b_n}
        + \Bigl(1 - \frac{b_N}{b_n}\Bigr)
        \Bigl(\frac{a_n - a_N}{b_n - b_N} - \ell\Bigr),
        \]
        并由此得出定理：$\dfrac{a_n}{b_n} \to \ell$。
  \item 证明 $+\infty$ 的变体：若 $\dfrac{a_{n+1} -
        a_n}{b_{n+1} - b_n} \to +\infty$（$(b_n)$ 的假设不变），
        则 $\dfrac{a_n}{b_n} \to +\infty$。
  \item 取 $b_n = n$，由此重新得到
        \cref{exo:b1:seq:10}中的切萨罗均值定理。再说明切萨罗--施托尔茨定理的逆命题不成立：
        对 $a_n = (-1)^n$，$b_n = n$，
        商 $a_n/b_n$ 收敛，而增量之商
        不收敛。施托尔茨定理只能单向使用。
\end{enumerate}

\medskip
\noindent\textbf{第三部分 --- 初步收获。}
\begin{enumerate}[resume]
  \item 证明 $(1+h)^{3/2} - 1 = \dfrac{3h + 3h^2 +
        h^3}{(1+h)^{3/2} + 1}$（用共轭法），由此推出
        $n\bigl((1 + \tfrac1n)^{3/2} - 1\bigr) \to \tfrac32$,
        并借切萨罗--施托尔茨定理得出：
        \[
        T_n = \sum_{k=1}^{n} \sqrt k \;\sim\; \tfrac23\,
        n^{3/2} ,
        \]
        从而解决第 2 题留下的悬念。
  \item 仅由（G1）推出对数的夹逼估计：
        \[
        \frac{t}{1 + t} \;\leq\; \ln(1 + t) \;\leq\; t
        \qquad (t > -1)
        \]
        \emph{（分别在 $u = \ln(1+t)$ 与 $u =
        -t/(1+t)$ 处应用（G1））}。
  \item 证明 $b_n = \ln n$ 严格递增，且
        $\ln n \to +\infty$，并借助切萨罗--施托尔茨定理与
        第 9 题证明下式：
        \[
        H_n = \sum_{k=1}^{n} \frac 1k \;\sim\; \ln n .
        \]
        （更精细的结构 $H_n = \ln n + \gamma + o(1)$ 是
        在\cref{ch:b1:series}中的周末问题。）
  \item （从比到根）设 $u_n > 0$，且
        $\frac{u_{n+1}}{u_n} \to L > 0$。利用第 9 题证明
        $\ln\frac{u_{n+1}}{u_n} \to \ln L$；对其应用切萨罗
        均值定理得 $\frac{\ln u_n}{n} \to \ln L$，再由（G1）
        得 $u_n^{1/n} \to L$。应用：计算
        $\lim\,\binom{2n}{n}^{1/n}$.
\end{enumerate}

\medskip
\noindent\textbf{第四部分 --- 正弦的缓慢下降。}设
        $u_0 \in \R$，且 $u_{n+1} = \sin u_n$。
\begin{enumerate}[resume]
  \item 由（G2）证明 $0 < \sin x < x$，其中 $0 < x \leq 1$。
        由此推出：$u_1 \in \intcc{-1}{1}$；若 $u_1 = 0$，则
        数列自第 $1$ 项起恒为零；若 $u_1 > 0$（
        $u_1 < 0$ 的情形对称，因 $\sin$ 为奇函数），则
        $(u_n)_{n \geq 1}$ 严格递减、为正，且
        收敛于 $0$ \emph{（由 $\ell =
        \sin \ell$ 求极限，并利用 $\abs{\sin a - \sin b} \leq \abs{a -
        b}$，此式本身由（G2）与和差化积公式推出
        ）}。
  \item 从现在起假设 $u_1 \in \intoc{0}{1}$，利用
        （G2），用夹逼法证明：
        \[
        \frac{\sin u_n}{u_n} \to 1
        \qquad\text{和}\qquad
        \frac{u_n - \sin u_n}{u_n^{3}} \to \frac16 .
        \]
  \item 证明如下分解式：
        \[
        w_n := \frac{1}{u_{n+1}^{2}} - \frac{1}{u_n^{2}}
        = \frac{u_n - \sin u_n}{u_n^{3}} \cdot
        \frac{u_n + \sin u_n}{u_n} \cdot
        \Bigl(\frac{u_n}{\sin u_n}\Bigr)^{2},
        \]
        并由此推出 $w_n \to \frac13$。
  \item 由\cref{exo:b1:seq:10}（增量形式）推出
        $\frac{1}{n\,u_n^{2}} \to \frac13$，再
        用平方根的共轭化方法，得主要结论：
        \[
        \sqrt n\;u_n \longrightarrow \sqrt 3 ,
        \qquad\text{即}\qquad
        u_n \sim \sqrt{\frac 3n} .
        \]
  \item 量化其缓慢程度：证明从某项起，$\sqrt{2/n}
        \leq u_n \leq 2/\sqrt n$，从而使 $u_n \leq
        10^{-2}$ 需要超过 $20\,000$ 次迭代（约
        $30\,000$ 次，由渐近式可得）。与海伦方法
        （见\cref{ex:b1:seq:heron}）比较，并说明其结构上的原因：
        即在不动点 $0$ 处，$\sin$ 的斜率
        为 $1$（\emph{中性}不动点），而每步误差减半的迭代
        需要斜率的模 $< 1$。
  \item 证明对\emph{任意}起点 $u_0 \in \R$，
        要么 $u_n = 0$（自第 $1$ 项起），要么 $\abs{u_n} \sim
        \sqrt{3/n}$ --- 下降是普遍的，只有符号
        记住了 $u_0$。
\end{enumerate}

\medskip
\noindent\textbf{第 V 部分 --- 一般原理。} 正弦只是
一台机器的一个实例。
\begin{enumerate}[resume]
  \item 设 $u_n > 0$，$u_n \to 0$，且 $\dfrac{u_n -
        u_{n+1}}{u_n^{2}} \to a > 0$。依次证明：
        $\frac{u_{n+1}}{u_n} \to 1$；再证 $\frac{1}{u_{n+1}} -
        \frac{1}{u_n} \to a$；再证 $n\,u_n \to \frac1a$。
  \item （精确模型）对 $u_{n+1} = \dfrac{u_n}{1 + u_n}$，$u_0
        > 0$：证明 $\frac{1}{u_n}$ 是等差数列，求出其通项，
        并将第 18 问的结论与
        精确公式比较。
  \item 对 $u_{n+1} = u_n \eu^{-u_n}$，$u_0 > 0$：证明 $u_n \to
        0$，利用 (G1) 将 $\frac{1 - \eu^{-t}}{t}$ 夹在
        $\frac{1}{1+t}$ 与 $1$ 之间（$t > 0$），并得出 $u_n
        \sim \frac 1n$。
  \item （三次接触，平方裂项）设 $u_n > 0$，$u_n
        \to 0$，$\dfrac{u_n - u_{n+1}}{u_n^{3}} \to a > 0$。
        仿照第 14 问的因式分解，证明
        $\frac{1}{u_{n+1}^2} - \frac{1}{u_n^2} \to 2a$，并
        得出 $n\,u_n^{2} \to \frac{1}{2a}$。检验 $a =
        \frac16$ 时可恢复第 IV 部分的结果。
\end{enumerate}

\medskip
\noindent\textbf{第 VI 部分 --- 方法的局限与启示。}
\begin{enumerate}[resume]
  \item 证明假设 $b_n \to +\infty$ 不可省去：
        对 $a_n = 2 - 2^{-n}$ 与 $b_n = 1 - 2^{-n}$，
        增量商趋于 $1$，而 $\frac{a_n}
        {b_n} \to 2$。指出第 5 问证明中
        失效的确切一行。
  \item （两次施托尔茨）证明 $\sum_{k=1}^{n} H_k \sim n \ln n$
        \emph{（一次应用切萨罗–施托尔茨，然后用第 10 问；
        利用第 9 问估计 $(n+1)\ln(n+1) - n\ln n$）}。
  \item （几何平均）若 $u_n > 0$ 且 $u_n \to \ell > 0$，
        证明 $(u_1 u_2 \cdots u_n)^{1/n} \to \ell$；若 $u_n \to
        +\infty$，证明 $(u_1 \cdots u_n)^{1/n} \to +\infty$。
        由此推出 $(n!)^{1/n} \to +\infty$。
  \item 综合题，每小题用一句话作答：(i) 本题中究竟在哪一步
        用到了完备性；(ii) 切萨罗–施托尔茨在何种意义上是
        离散的洛必达法则（其微分版本
        依赖于
        \cref{ch:b1:derivative}中的中值定理）；(iii) 写出将
        $f$ 在中性不动点处的接触阶与
        $u_{n+1} = f(u_n)$ 的衰减指数相联系的启发式；
        (iv) 追溯 $3$ 在 $\sqrt{3/n}$ 中的来历，沿
        流程 $\frac16 \to \frac13 \to 3$ 逐步回溯。
\end{enumerate}
\end{problem}
